%% ma: L10-B14-0045
%% lop: 10
%% chuong: 5
%% bai: 14
%% dang: 10.14.9
%% loai: TL
%% muc: VD
%% dap_an: "-"
%% nguon: "chuyen-de-cac-so-dac-trung-cua-mau-so-lieu-khong-ghep-nhom-toan-10-chuong-trinh-moi.pdf (GV. Phan Nhật Linh) (P05), Bài 03 Dạng 2 Bài tập tự luận 5"
%% kien_thuc: "Lớp 10 Bài 13 (số trung bình), Bài 14 (phương sai, độ lệch chuẩn)"
%% trang_thai: chinh-thuc
%% ghi_chu: "Tự duyệt theo yêu cầu giáo viên 22/10/2026."
\begin{ex}
Dưới đây là điểm kiểm tra của hai bạn Ngân và Hà:
\begin{center}
\begin{tabular}{|c|c|c|c|c|c|c|c|c|c|c|c|}\hline
 & Toán & Lí & Hóa & Sinh & Anh & Văn & Sử & Địa & GDCD & Tin & CN\\\hline
Ngân & $8$ & $7{,}5$ & $7{,}8$ & $8{,}3$ & $9$ & $7$ & $8$ & $8{,}2$ & $9$ & $8$ & $8{,}3$\\\hline
Hà & $8{,}5$ & $9{,}5$ & $9{,}5$ & $8{,}5$ & $9$ & $5$ & $5{,}5$ & $6$ & $10$ & $9$ & $8{,}5$\\\hline
\end{tabular}
\end{center}
Hãy tính điểm trung bình các môn học của Ngân và Hà. Theo em, bạn nào học đều hơn?
\loigiai{$\overline x_N=\dfrac{89{,}1}{11}=8{,}1$; $\overline x_H=\dfrac{89}{11}\approx8{,}09$ (xấp xỉ bằng nhau).

$s_N^2=\dfrac{1}{11}\sum\left(x_i-8{,}1\right)^2=\dfrac{3{,}4}{11}\approx0{,}31$; $s_H^2\approx2{,}76$.

Điểm trung bình gần như bằng nhau và $s_N^2<s_H^2$ nên Ngân học đều hơn.}
\end{ex}
